\documentclass[a4paper,12pt]{article}
\usepackage[utf8]{inputenc}
\usepackage[T1]{fontenc}
\usepackage[ngerman]{babel}
\usepackage{amsmath, amssymb}
\usepackage{geometry}
\usepackage{tikz}
\usetikzlibrary{shapes.geometric, arrows.meta, positioning, fit, backgrounds}
\usepackage{parskip}
\usepackage{titlesec}
\usepackage{enumitem}

\geometry{left=3cm, right=3cm, top=2.5cm, bottom=2.5cm}

\titleformat{\section}{\large\bfseries}{}{0em}{}[\titlerule]
\titleformat{\subsection}{\normalsize\bfseries}{}{0em}{}

\tikzset{
  box/.style={rectangle, draw, rounded corners=3pt, text width=4.5cm,
              align=center, minimum height=1.0cm, font=\small, inner sep=4pt},
  boxw/.style={box, text width=5.5cm},
  arrow/.style={-{Latex[length=2mm]}, thick},
}

\begin{document}

\begin{center}
  {\LARGE\bfseries Studierhinweise zu Lektion 3}\\[0.5em]
  {\large Mathematische Grundlagen (Modul 61111)}
\end{center}

\vspace{1em}

Diese Lektion bildet das Herzstück der Linearen Algebra im Modul. In
Kapitel~7 werden die Begriffe der linearen Unabhängigkeit, der Basis und der
Dimension endlich erzeugter Vektorräume eingeführt. Das zentrale Hilfsmittel
ist der \emph{Austauschsatz von Steinitz}, aus dem alle wichtigen
Dimensionsformeln folgen.

Kapitel~8 wendet sich den \emph{linearen Abbildungen} zu – Abbildungen
zwischen Vektorräumen, die Addition und Skalarmultiplikation respektieren.
Hier lernen Sie die Begriffe Kern, Bild und Rang einer linearen Abbildung
kennen sowie den wichtigen \emph{Rangsatz}.

Kapitel~9 schlägt die Brücke zurück zu Matrizen: Jeder linearen Abbildung
zwischen endlich erzeugten Vektorräumen kann (nach Wahl von Basen) eine
Matrix zugeordnet werden. Matrizenmultiplikation entspricht dabei der
Komposition linearer Abbildungen.

Lernen Sie die Definitions- und Satzformulierungen so weit, dass Sie sie
erkennen, anwenden und — zumindest skizzenhaft — beweisen können. Viele
Sätze dieser Lektion sind Folgerungen aus dem Steinitz'schen Austauschsatz
und bauen aufeinander auf.

% -----------------------------------------------------------------------
\section*{Struktur der Lektion 3}

\begin{center}
\begin{tikzpicture}[node distance=0.55cm and 1.0cm]

  % Headers
  \node[font=\bfseries\small] (h7) at (0,0)    {Kapitel 7: Basen \& Dimension};
  \node[font=\bfseries\small] (h8) at (6.5,0)  {Kapitel 8: Lineare Abbildungen};
  \node[font=\bfseries\small] (h9) at (12.5,0) {Kapitel 9: Abbildungen \& Matrizen};

  % Kapitel 7
  \node[box, below=0.8cm of h7] (lu)   {7.1 Lineare Unabhängigkeit\\trivial / nicht-trivial};
  \node[box, below=0.5cm of lu]  (bas) {7.2 Basen\\Charakterisierung, Existenz};
  \node[box, below=0.5cm of bas] (aus) {7.3 Austauschsatz (Steinitz)\\Basisergänzungssatz};
  \node[box, below=0.5cm of aus] (dim) {7.4 Dimension\\Dimensionsformeln};

  % Kapitel 8
  \node[box, below=0.8cm of h8] (def)  {8.1 Definition \& Eigenschaften\\$f(v_1+v_2)=f(v_1)+f(v_2)$};
  \node[box, below=0.5cm of def] (iso) {8.2 Isomorphe Vektorräume\\$\kappa_B : V \to K^n$};
  \node[box, below=0.5cm of iso] (kb)  {8.3 Kern \& Bild\\Rangsatz};
  \node[box, below=0.5cm of kb]  (ba)  {8.4 Abbildungen \& Basen\\Eindeutigkeit durch Bilder};

  % Kapitel 9
  \node[box, below=0.8cm of h9] (hom) {9.1 $\mathrm{Hom}_K(V,W)$\\Vektorraumstruktur};
  \node[box, below=0.5cm of hom] (koo) {9.2 Koordinatenvektoren\\$_C M_B(f)\cdot\kappa_B(v)=\kappa_C(f(v))$};
  \node[box, below=0.5cm of koo] (kom) {9.3 Matrizenprodukt \&\\Komposition};

  % Arrows Kap 7
  \draw[arrow] (lu)  -- (bas);
  \draw[arrow] (bas) -- (aus);
  \draw[arrow] (aus) -- (dim);

  % Arrows Kap 8
  \draw[arrow] (def) -- (iso);
  \draw[arrow] (iso) -- (kb);
  \draw[arrow] (kb)  -- (ba);

  % Arrows Kap 9
  \draw[arrow] (hom) -- (koo);
  \draw[arrow] (koo) -- (kom);

  % Cross-chapter
  \draw[arrow] (dim.east) -- ++(0.3,0) |- (def.west);
  \draw[arrow] (ba.east)  -- ++(0.3,0) |- (hom.west);

\end{tikzpicture}
\end{center}

\newpage

% -----------------------------------------------------------------------
\section*{Zielelement 7.1 -- Lineare Unabhängigkeit}

\subsection*{Lerninhalte}

\begin{center}
\begin{tikzpicture}[node distance=0.55cm and 1.0cm]
  \node[box] (triv) {Triviale Darstellung des Nullvektors (alle $a_i = 0$)};
  \node[box, below=0.5cm of triv] (def) {Definition: linear unabhängig\\$\sum a_i v_i = 0 \Rightarrow a_i = 0 \;\forall i$};
  \node[box, below=0.5cm of def]  (lab) {Linear abhängig:\\nicht-triviale Darstellung von $0$};
  \node[box, below=0.5cm of lab]  (cha) {Prop.~7.1.6: Charakterisierung\\$v_i \in \langle v_1,\ldots,\hat{v}_i,\ldots,v_n\rangle$};
  \node[box, below=0.5cm of cha]  (gau) {Kriterium: Rang der Koeffizientenmatrix};
  \draw[arrow] (triv) -- (def);
  \draw[arrow] (def)  -- (lab);
  \draw[arrow] (lab)  -- (cha);
  \draw[arrow] (cha)  -- (gau);
\end{tikzpicture}
\end{center}

\subsection*{Lernziele}

Nach Durcharbeiten dieses Abschnitts sollten Sie:
\begin{itemize}
  \item den Begriff der linearen Unabhängigkeit präzise definieren können,
  \item lineare Unabhängigkeit von Vektoren mithilfe des Gauß-Algorithmus
        (Rangberechnung) überprüfen können,
  \item die Charakterisierung linear abhängiger Vektoren (Prop.~7.1.6) kennen
        und anwenden können,
  \item wissen, dass ein einzelner Vektor $v \neq 0$ stets linear unabhängig ist, und
        zwei Vektoren genau dann linear abhängig sind, wenn einer ein skalares
        Vielfaches des anderen ist.
\end{itemize}

\subsection*{Selbstkontrollelement 7.1}

Sind die Vektoren $\begin{pmatrix}1\\2\\0\end{pmatrix}$,
$\begin{pmatrix}0\\1\\3\end{pmatrix}$,
$\begin{pmatrix}2\\3\\-3\end{pmatrix} \in \mathbb{R}^3$ linear unabhängig?
Bestimmen Sie dazu den Rang der zugehörigen Koeffizientenmatrix.

\newpage

% -----------------------------------------------------------------------
\section*{Zielelement 7.2 -- Basen endlich erzeugter Vektorräume}

\subsection*{Lerninhalte}

\begin{center}
\begin{tikzpicture}[node distance=0.55cm and 1.0cm]
  \node[box] (def) {Def.~7.2.1: Basis\\linear unabhängiges Erzeugendensystem};
  \node[box, below=0.5cm of def]  (std) {Standardbasen: $K^n$, $M_{mn}(K)$,\\Polynome vom Grad $\le n$};
  \node[box, below=0.5cm of std]  (cha) {Prop.~7.2.7: Äquivalenz\\eindeutige Darstellung jedes $v \in V$};
  \node[box, below=0.5cm of cha]  (exi) {Prop.~7.2.8: Existenz von Basen\\Herausstreichen aus Erzeugendensystem};
  \node[box, below=0.5cm of exi]  (bem) {Bem.~7.2.11: linear unabh.\ Vektoren\\ergänzen / Erzeugendensystem};
  \draw[arrow] (def) -- (std);
  \draw[arrow] (std) -- (cha);
  \draw[arrow] (cha) -- (exi);
  \draw[arrow] (exi) -- (bem);
\end{tikzpicture}
\end{center}

\subsection*{Lernziele}

Nach Durcharbeiten dieses Abschnitts sollten Sie:
\begin{itemize}
  \item die Definition einer Basis auswendig kennen,
  \item die Standardbasen der wichtigsten Vektorräume ($K^n$, $M_{mn}(K)$,
        Polynome) angeben können,
  \item Proposition~7.2.7 (Charakterisierung von Basen durch eindeutige
        Linearkombinationsdarstellung) kennen und im Beweis nutzen können,
  \item verstehen, wie man aus einem endlichen Erzeugendensystem eine Basis
        gewinnt (Proposition~7.2.8).
\end{itemize}

\subsection*{Selbstkontrollelement 7.2}

Zeigen Sie, dass $1+T$, $1-T$ eine Basis des Vektorraums der Polynome vom
Grad $\le 1$ über $\mathbb{R}$ ist. Warum ist $1+T$, $1-T$, $2$ keine Basis?

\newpage

% -----------------------------------------------------------------------
\section*{Zielelement 7.3 -- Der Austauschsatz von Steinitz}

\subsection*{Lerninhalte}

\begin{center}
\begin{tikzpicture}[node distance=0.55cm and 1.0cm]
  \node[box] (lem) {Lemma~7.3.1: Austauschlemma\\Ersatz eines Basisvektors durch $v \ne 0$};
  \node[box, below=0.5cm of lem] (satz){Satz~7.3.3: Austauschsatz (Steinitz)\\$m$ lin.\ unabh.\ Vektoren $\Rightarrow m \le n$};
  \node[box, below=0.5cm of satz] (kor1){Kor.~7.3.4: Endlich erzeugt $\Leftrightarrow$\\beschränkte Anzahl lin.\ unabh.\ Vektoren};
  \node[box, below=0.5cm of kor1] (kor2){Kor.~7.3.5: Basisergänzungssatz\\lin.\ unabh.\ Menge zu Basis ergänzen};
  \node[box, below=0.5cm of kor2] (kor3){Kor.~7.3.6: Alle Basen\\haben gleiche Elementezahl};
  \draw[arrow] (lem)  -- (satz);
  \draw[arrow] (satz) -- (kor1);
  \draw[arrow] (kor1) -- (kor2);
  \draw[arrow] (kor2) -- (kor3);
\end{tikzpicture}
\end{center}

\subsection*{Lernziele}

Nach Durcharbeiten dieses Abschnitts sollten Sie:
\begin{itemize}
  \item das Austauschlemma (7.3.1) und den Austauschsatz (7.3.3) in eigenen
        Worten erklären können,
  \item den Basisergänzungssatz (7.3.5) kennen und anwenden können,
  \item begründen können, warum je zwei Basen eines endlich erzeugten
        Vektorraums gleich viele Elemente haben (7.3.6).
\end{itemize}

\subsection*{Selbstkontrollelement 7.3}

Sei $V = \mathbb{R}^3$. Ergänzen Sie $v_1 = \begin{pmatrix}1\\1\\0\end{pmatrix}$
zu einer Basis von $V$. Wie viele Vektoren müssen Sie hinzufügen?

\newpage

% -----------------------------------------------------------------------
\section*{Zielelement 7.4 -- Dimension und Dimensionsformeln}

\subsection*{Lerninhalte}

\begin{center}
\begin{tikzpicture}[node distance=0.55cm and 1.0cm]
  \node[box] (def) {Def.~7.4.1: Dimension $\dim(V)$\\Anzahl Vektoren einer Basis};
  \node[box, below=0.5cm of def]  (bsp) {Beispiele: $\dim(K^n)=n$,\\$\dim(M_{mn})=mn$, $\dim(\text{Pol}_n)=n+1$};
  \node[box, below=0.5cm of bsp]  (pro) {Prop.~7.4.4: $n$ lin.\ unabh.\ Vektoren\\in $V$ mit $\dim(V)=n$ $\Rightarrow$ Basis};
  \node[box, below=0.5cm of pro]  (unt) {Kor.~7.4.5: $\dim(U) \le \dim(V)$\\$\dim(U)=\dim(V) \Leftrightarrow U=V$};
  \node[box, below=0.5cm of unt]  (sum) {Prop.~7.4.6: Dimensionsformel\\$\dim(U+W)=\dim(U)+\dim(W)-\dim(U\cap W)$};
  \draw[arrow] (def) -- (bsp);
  \draw[arrow] (bsp) -- (pro);
  \draw[arrow] (pro) -- (unt);
  \draw[arrow] (unt) -- (sum);
\end{tikzpicture}
\end{center}

\subsection*{Lernziele}

Nach Durcharbeiten dieses Abschnitts sollten Sie:
\begin{itemize}
  \item die Dimension eines Vektorraums definieren und für die Standardbeispiele
        angeben können,
  \item Proposition~7.4.4 kennen und als Arbeitserleichterung einsetzen können
        (nur eine der beiden Basiseigenschaften nachweisen),
  \item die Dimensionsformel für Summe und Durchschnitt (Prop.~7.4.6) beherrschen
        und den Beweisaufbau skizzieren können.
\end{itemize}

\subsection*{Selbstkontrollelement 7.4}

Seien $U = \langle e_1, e_2 \rangle$ und $W = \langle e_2, e_3 \rangle$ Unterräume
von $\mathbb{R}^3$. Bestimmen Sie $\dim(U+W)$ und $\dim(U \cap W)$
mithilfe der Dimensionsformel.

\newpage

% -----------------------------------------------------------------------
\section*{Zielelement 8.1 -- Lineare Abbildungen: Definition und Eigenschaften}

\subsection*{Lerninhalte}

\begin{center}
\begin{tikzpicture}[node distance=0.55cm and 1.0cm]
  \node[box] (def) {Def.~8.1.1: Lineare Abbildung\\$f(v_1+v_2)=f(v_1)+f(v_2)$,\\$f(av)=af(v)$};
  \node[box, below=0.5cm of def]  (bem) {Bem.~8.1.4: $f(0)=0$,\\$f(-v)=-f(v)$};
  \node[box, below=0.5cm of bem]  (iso) {Def.~8.1.5: Isomorphismus\\bijektive lineare Abbildung};
  \node[box, below=0.5cm of iso]  (inv) {Prop.~8.1.8: Inverse von\\Isomorphismen sind Isomorphismen};
  \node[box, below=0.5cm of inv]  (kom) {Prop.~8.1.9: Komposition\\linearer Abbildungen ist linear};
  \draw[arrow] (def) -- (bem);
  \draw[arrow] (bem) -- (iso);
  \draw[arrow] (iso) -- (inv);
  \draw[arrow] (inv) -- (kom);
\end{tikzpicture}
\end{center}

\subsection*{Lernziele}

Nach Durcharbeiten dieses Abschnitts sollten Sie:
\begin{itemize}
  \item lineare Abbildungen erkennen und die Linearitätsbedingungen nachweisen
        können,
  \item die elementaren Eigenschaften linearer Abbildungen (Nullvektor, negierter Vektor)
        aus den Axiomen herleiten können,
  \item den Begriff des Isomorphismus kennen und verstehen, warum inverse und
        zusammengesetzte Isomorphismen wieder Isomorphismen sind.
\end{itemize}

\subsection*{Selbstkontrollelement 8.1}

Ist $f: \mathbb{R}^2 \to \mathbb{R}^2$, $f\!\begin{pmatrix}x\\y\end{pmatrix} = \begin{pmatrix}x+1\\y\end{pmatrix}$
eine lineare Abbildung? Begründen Sie.

\newpage

% -----------------------------------------------------------------------
\section*{Zielelement 8.2 -- Isomorphe Vektorräume}

\subsection*{Lerninhalte}

\begin{center}
\begin{tikzpicture}[node distance=0.55cm and 1.0cm]
  \node[box] (kb)  {Satz~8.2.2: $\kappa_B: V \to K^n$\\ist Isomorphismus};
  \node[box, below=0.5cm of kb]  (koo) {Def.~8.2.4: Koordinatenvektor\\$\kappa_B(v) \in K^n$ bzgl.\ Basis $B$};
  \node[box, below=0.5cm of koo] (kor) {Kor.~8.2.6:\\gleiche Dim.\ $\Leftrightarrow$ isomorph};
  \node[box, below=0.5cm of kor] (str) {Satz~8.3.19: Struktursatz\\$\dim(V)=n \Rightarrow V \cong K^n$};
  \draw[arrow] (kb)  -- (koo);
  \draw[arrow] (koo) -- (kor);
  \draw[arrow] (kor) -- (str);
\end{tikzpicture}
\end{center}

\subsection*{Lernziele}

Nach Durcharbeiten dieses Abschnitts sollten Sie:
\begin{itemize}
  \item den Koordinatenvektor eines Vektors bezüglich einer Basis bestimmen können,
  \item die Abbildung $\kappa_B$ als Isomorphismus verstehen und erklären können,
  \item wissen, dass zwei endlich erzeugte Vektorräume genau dann isomorph sind,
        wenn sie dieselbe Dimension haben.
\end{itemize}

\subsection*{Selbstkontrollelement 8.2}

Sei $B = (1+T,\, T)$ eine Basis des Vektorraums der Polynome vom Grad $\le 1$
über $\mathbb{R}$. Bestimmen Sie $\kappa_B(3 + 2T)$.

\newpage

% -----------------------------------------------------------------------
\section*{Zielelement 8.3 -- Kern, Bild und Rangsatz}

\subsection*{Lerninhalte}

\begin{center}
\begin{tikzpicture}[node distance=0.55cm and 1.0cm]
  \node[box] (bil) {Def.~8.3.1: Bild$(f)$\\Unterraum von $W$ (Bem.~8.3.2)};
  \node[box, right=1.0cm of bil] (ker) {Def.~8.3.7: Kern$(f)$\\Unterraum von $V$ (Bem.~8.3.8)};
  \node[box, below=0.5cm of bil]  (sur) {Kor.~8.3.5: $f$ surjektiv\\$\Leftrightarrow f(v_1),\ldots,f(v_n)$ Erz.sys.\ von $W$};
  \node[box, below=0.5cm of ker]  (inj) {Prop.~8.3.10: $f$ injektiv\\$\Leftrightarrow \mathrm{Kern}(f) = \{0\}$};
  \node[box, below=0.5cm of sur, xshift=2.5cm] (rg)  {Def.~8.3.12: $\mathrm{Rg}(f) = \dim(\mathrm{Bild}(f))$};
  \node[box, below=0.5cm of rg]  (rs)  {Kor.~8.3.14: Rangsatz\\$\dim(V) = \dim(\mathrm{Kern}(f)) + \mathrm{Rg}(f)$};
  \node[box, below=0.5cm of rs]  (kor) {Kor.~8.3.17: surjektiv $\Leftrightarrow$ injektiv\\(bei $\dim(V)=\dim(W)$)};
  \draw[arrow] (bil) -- (sur);
  \draw[arrow] (ker) -- (inj);
  \draw[arrow] (sur.south) -- ++(0,-0.2) -| (rg.north);
  \draw[arrow] (inj.south) -- ++(0,-0.2) -| (rg.north);
  \draw[arrow] (rg)  -- (rs);
  \draw[arrow] (rs)  -- (kor);
\end{tikzpicture}
\end{center}

\subsection*{Lernziele}

Nach Durcharbeiten dieses Abschnitts sollten Sie:
\begin{itemize}
  \item Kern und Bild einer linearen Abbildung berechnen können,
  \item Injektivität über den Kern und Surjektivität über das Bild charakterisieren können,
  \item den Rangsatz kennen, beweisen und auf konkrete Beispiele anwenden können,
  \item Korollar~8.3.17 (surjektiv $\Leftrightarrow$ injektiv) als Arbeitserleichterung nutzen.
\end{itemize}

\subsection*{Selbstkontrollelement 8.3}

Sei $f: \mathbb{R}^3 \to \mathbb{R}^2$, $f\!\begin{pmatrix}x\\y\\z\end{pmatrix} = \begin{pmatrix}x+y\\y+z\end{pmatrix}$.
Bestimmen Sie $\mathrm{Kern}(f)$ und $\mathrm{Rg}(f)$. Verifizieren Sie den Rangsatz.

\newpage

% -----------------------------------------------------------------------
\section*{Zielelement 8.4 -- Lineare Abbildungen und Basen}

\subsection*{Lerninhalte}

\begin{center}
\begin{tikzpicture}[node distance=0.55cm and 1.0cm]
  \node[box] (satz) {Satz~8.4.1: Lineare Abbildung\\eindeutig durch Bilder einer Basis};
  \node[box, below=0.5cm of satz] (kons){Konstruktion linearer Abbildungen\\mit vorgegebenen Eigenschaften};
  \node[box, below=0.5cm of kons] (kor1){Kor.~8.4.6: Zu jeder Basis $B$, $C$\\gibt es genau einen Iso.\ $f(v_i)=w_i$};
  \node[box, below=0.5cm of kor1] (kor2){Kor.~8.4.7 / Def.~8.4.8:\\Basiswechsel als Isomorphismus};
  \draw[arrow] (satz) -- (kons);
  \draw[arrow] (kons) -- (kor1);
  \draw[arrow] (kor1) -- (kor2);
\end{tikzpicture}
\end{center}

\subsection*{Lernziele}

Nach Durcharbeiten dieses Abschnitts sollten Sie:
\begin{itemize}
  \item Satz~8.4.1 kennen und zur Konstruktion linearer Abbildungen mit
        vorgegebenen Eigenschaften einsetzen können,
  \item verstehen, warum eine lineare Abbildung vollständig durch ihre Werte
        auf einer Basis festgelegt ist,
  \item den Begriff des Basiswechsels (Def.~8.4.8) erklären können.
\end{itemize}

\subsection*{Selbstkontrollelement 8.4}

Definieren Sie eine lineare Abbildung $f: \mathbb{R}^2 \to \mathbb{R}^2$ mit
$\mathrm{Kern}(f) = \langle\begin{pmatrix}1\\1\end{pmatrix}\rangle$.
Geben Sie eine explizite Formel für $f$ an.

\newpage

% -----------------------------------------------------------------------
\section*{Zielelement 9 -- Lineare Abbildungen und Matrizen}

\subsection*{Lerninhalte}

\begin{center}
\begin{tikzpicture}[node distance=0.55cm and 1.0cm]
  \node[box] (hom) {Def.~9.1.1: $\mathrm{Hom}_K(V,W)$\\Vektorraumstruktur (Prop.~9.1.2)};
  \node[box, below=0.5cm of hom] (mat) {Def.~9.1.4: Matrixdarstellung $_C M_B(f)$\\Koordinatenvektoren der Bilder};
  \node[box, below=0.5cm of mat] (iso) {Satz~9.1.9: $_C M_B$ ist Isomorphismus\\$\mathrm{Hom}_K(V,W) \cong M_{mn}(K)$};
  \node[box, below=0.5cm of iso] (dim) {Kor.~9.1.10:\\$\dim(\mathrm{Hom}_K(V,W)) = \dim(V)\cdot\dim(W)$};
  \node[box, below=0.5cm of dim] (koo) {Prop.~9.2.1: $_C M_B(f)\cdot\kappa_B(v) = \kappa_C(f(v))$};
  \node[box, below=0.5cm of koo] (rg)  {Kor.~9.2.4: $\mathrm{Rg}(f) = \mathrm{Rg}(_C M_B(f))$};
  \node[box, below=0.5cm of rg]  (pro) {Prop.~9.3.1: $_D M_B(g\circ f) = {}_D M_C(g)\cdot{}_C M_B(f)$};
  \draw[arrow] (hom) -- (mat);
  \draw[arrow] (mat) -- (iso);
  \draw[arrow] (iso) -- (dim);
  \draw[arrow] (dim) -- (koo);
  \draw[arrow] (koo) -- (rg);
  \draw[arrow] (rg)  -- (pro);
\end{tikzpicture}
\end{center}

\subsection*{Lernziele}

Nach Durcharbeiten dieses Kapitels sollten Sie:
\begin{itemize}
  \item die Matrixdarstellung einer linearen Abbildung bezüglich gegebener Basen
        berechnen können,
  \item die zentrale Formel $_{C}M_B(f)\cdot\kappa_B(v) = \kappa_C(f(v))$ verstehen
        und anwenden können,
  \item wissen, dass Matrizenmultiplikation der Komposition linearer Abbildungen
        entspricht (Prop.~9.3.1),
  \item die Dimensionsformel $\dim(\mathrm{Hom}_K(V,W)) = mn$ kennen.
\end{itemize}

\subsection*{Selbstkontrollelement 9}

Sei $f: \mathbb{R}^2 \to \mathbb{R}^2$ definiert durch $f\!\begin{pmatrix}x\\y\end{pmatrix} = \begin{pmatrix}2x-y\\x+3y\end{pmatrix}$.
Berechnen Sie $_E M_E(f)$ bezüglich der Standardbasis $E$ von $\mathbb{R}^2$.
Was ist $\mathrm{Rg}(f)$?

\end{document}
